% ===== 原文 pp.28--29（PDF pp.29--30）=====
% tr-03 片段 c：$\mathrm{T}^{U(1)}_{N,d}(0)$ 模型、$\mathrm{H}_n$ 模型开头（到 (4.34)）、Figure 3

% [原文 p.28]
\para{模型 $\mathrm{T}^{U(1)}_{N,d}(0)$}

我们也可以考虑模型 $\mathrm{T}^{U(1)}_{N,d}(0)$：它与上一例有相同的 $(G,V,\tW)$，但超势取零 $W=0$。对矢量 R-荷的约束（即 $W$ 必须具有 R-荷 2）不再存在，结果最自然的选择是令所有场 $p,x_1,\ldots,x_N$ 的 R-荷都为 $0$。

标量势缺少 (4.29) 中的最后两项，相应地，经典相的描述也随之改变。

当 $\zeta>0$ 时，$x$ 仍非零并完全破缺规范群，从而 $\sigma=0$。不过这一次 $p$ 不必为零。真空流形现在是满足 $-d|p|^2+|x|^2=\zeta$ 的 $(p,x)$ 的空间模去规范 $U(1)$ 作用。它是 $\CP^{N-1}$ 上全纯线丛 $\cO_{\CP^{N-1}}(-d)$ 的全空间，其中 $p$ 与 $x$ 分别参数化纤维与底。所有横截模式都是有质量的，理论在低能下约化为靶空间为非紧的 $\mathrm{tot}[\cO_{\CP^{N-1}}(-d)]$ 的非线性 sigma 模型。

当 $\zeta<0$ 时，$p$ 非零并把规范群破缺到 $\Z_d$，从而 $\sigma=0$。只有 $x$ 场是无质量的，于是我们剩下 $G=\Z_d$、$V=\C(1)^{\oplus N}$、$W=0$（而 $\tW=0$，理由平凡）的理论。这就是自由轨形理论 $\C^N/\Z_d$。注意 $\mathrm{tot}[\cO_{\CP^{N-1}}(-d)]$ 可以看作原点处轨形奇点的一个消解。

当 $\zeta=0$ 时，真空流形是两个分支的并——$\sigma=0$ 的 Higgs 分支与 $x=p=0$ 的 Coulomb 分支。它们在原点 $x=p=\sigma=0$ 处彼此相触。

在量子理论中，FI 参数的跑动、轴向异常，以及大 $\sigma$ 值处的有效扭超势，都与有超势的模型相同。特别地，Coulomb 真空的方程仍与 (4.32) 相同。细节依赖于 $N$ 与 $d$ 的关系。

当 $d=N$ 时，我们得到一族由 $\ee^t$ 参数化的非紧超共形场论。判别轨迹为 $\ee^t=(-N)^N$。我们看到，靶为 $\mathrm{tot}[\cO_{\CP^{N-1}}(-N)]$ 的 sigma 模型与自由轨形 $\C^N/\Z_N$ 连续相连。特别地，二者的拓扑 B-模型是等价的，而拓扑 A-模型由解析延拓相联系。$N=2$ 的情形是 \emph{McKay 对应}的一个例子，其余情形则是它的推广。

当 $d<N$ 时，FI 参数 $\zeta$ 由正跑到负，使得 $\Lambda^{N-d}:=\ee^{-t}\mu^{N-d}$ 是 RG 不变量。高能下的理论是靶为 $\mathrm{tot}[\cO_{\CP^{N-1}}(-d)]$ 的非线性 sigma 模型。存在 $(N-d)$ 个有质量间隙的 Coulomb 真空，位于 $\sigma^{N-d}=(-d)^d\Lambda^{N-d}$；以及在 $\sigma=0$ 处的一个 Higgs 分支理论（对 $d>1$），它就是
% [原文 p.29]
自由轨形 $\C^N/\Z_d$。

当 $d>N$ 时，FI 参数由负跑到正，使得 $\Lambda^{d-N}:=\ee^{t}\mu^{d-N}$ 是 RG 不变量。该理论可视为自由轨形 $\C^N/\Z_d$ 的相关形变。存在 $(d-N)$ 个有质量间隙的 Coulomb 真空，位于 $\sigma^{d-N}=(-d)^{-d}\Lambda^{d-N}$；以及在 $\sigma=0$ 处的一个 Higgs 分支理论，它是靶为 $\mathrm{tot}[\cO_{\CP^{N-1}}(-d)]$ 的非线性 sigma 模型。

\para{模型 $\mathrm{H}_n$}

最后，考虑由非负整数 $n$ 标记的模型 $\mathrm{H}_n$：
\begin{align*}
    G &= U(1)\times U(1),\\
    V &= \C(1,0)\oplus\C(1,0)\oplus\C(0,1)\oplus\C(-n,1)\ni(\phi_1,\phi_2,\phi_3,\phi_4),\\
    W &= 0,\\
    \tW &= -t^1\sigma_1-t^2\sigma_2.
\end{align*}

\begin{figure}[htbp]
\centering
\includegraphics[width=0.62\textwidth]{fig3.png}
\figtitle{3}{$\mathrm{H}_n$ 的各相（$n=3$）。}
\end{figure}

写出标量势的工作留给读者作为练习。经典超对称真空方程由 $\sigma_1\phi_1=\sigma_1\phi_2=\sigma_2\phi_3=(-n\sigma_1+\sigma_2)\phi_4=0$ 以及
\begin{equation}
    |\phi_1|^2+|\phi_2|^2-n|\phi_4|^2=\zeta^1,\qquad
    |\phi_3|^2+|\phi_4|^2=\zeta^2. \tag{4.34}
\end{equation}
组成。
% [Tutor Note: p.29 式 (4.34) 就是两个 $U(1)$ 的 D-项（矩映射）方程：各场的规范荷为 $\phi_1,\phi_2:(1,0)$、$\phi_3:(0,1)$、$\phi_4:(-n,1)$，逐分量写出 $\sum_aQ_a|\phi_a|^2=\zeta$ 即得 (4.34)；这也说明 p.29 图 3 中各相由 $\zeta=(\zeta^1,\zeta^2)$ 平面上的腔室决定（图中标注的 $\phi_{1,2}$、$\phi_3$、$\phi_4$ 分别是各相中必须非零的场）。另一条可核对的线索：令 $\phi_4=0$，则 (4.34) 退化为 $|\phi_1|^2+|\phi_2|^2=\zeta^1$、$|\phi_3|^2=\zeta^2$，要它有解必须 $\zeta^1>0$、$\zeta^2>0$，即图 3 的 Phase I（第一象限）。]
% [Tutor Note: 原文笔误（正文照译、未改，均以 200 dpi 页面图核对过字面）：p.27 "corrsponding" 应为 "corresponding"；p.28 "at large σ0" 中的 "0" 疑为衍字（同段语境只涉及大 $\sigma$，即 (4.31) 之前一句），"most natural choise" 应为 "choice"；p.29 "non-negtaive" 应为 "non-negative"，"topologival A-models" 应为 "topological A-models"。]
